\begin{table}[!htbp]
\centering
\caption{Expected Consumption-Equivalent Welfare Changes from Tax Reform}
\label{tab:ce_welfare_tax}
\begin{threeparttable}
\begin{tabular}{lcc}
\toprule
& \textbf{Infinite horizon} & \textbf{Life-cycle} \\
\midrule
WT vs CIT       & 0.20\% & 0.15\% \\
WT vs No Tax  & 0.85\% & 0.35\% \\
CIT vs No Tax & 0.65\% & 0.20\% \\
\bottomrule
\end{tabular}
\begin{tablenotes}[flushleft]
\footnotesize
\item Notes: Entries are consumption-equivalent (CE) welfare changes, $\Delta$, expressed as percent changes. For example, WT vs CIT is the percentage increase in consumption under the wealth tax required to make a newborn indifferent to the capital income tax. Positive values therefore indicate that the right-hand policy yields higher newborn lifetime welfare than the left-hand policy.
\end{tablenotes}
\end{threeparttable}
\end{table}
